\section{Formalizing Handles and Their Guarantees}
\label{sec:model}

This section gives a compact model of the handle lifecycle and states the four guarantees claimed so far. The results are conditional on explicit reference-monitor assumptions; they characterize the model, not a deployed implementation, and are validated only on controlled synthetic inputs (Section~\ref{sec:evidence}).

\paragraph{Records and handles.} The trusted store maps each identifier $n$ to a record
\begin{equation}
\Omega[n]=(s,H,a,z,r),\qquad H=(n,q,\sigma,d,\tau,t_i,t_e,u,c),
\label{eq:record}
\end{equation}
where $s$ is the value; $q,\sigma,\tau$ are the type, provenance, and task identifier; $d$ is the destination predicate; $t_i,t_e$ are the issue and expiry times; $u$ is the remaining-use budget; $c$ is the confirmation policy; $a=F_K(\enc(H))$ is an authenticator over a canonical encoding of $H$; $z\in\{\mathsf{live},\mathsf{reserved},\mathsf{consumed},\mathsf{expired},\mathsf{revoked}\}$ is the lifecycle state; and $r$ is an optional reservation nonce. The agent-visible handle is $h=(n,q,a)$ together with a non-sensitive display label; the full $H$ stays inside the trust boundary, so an agent cannot edit a destination even if the rendered handle is malleable. Because the running use count changes over a handle's life while $\enc(H)$ must stay fixed for $a$ to remain valid, the store keeps $H$ immutable and tracks the remaining uses as mutable record state: the reservation step decrements it and finalization settles only the terminal state.

\begin{definition}[Target match]
\label{def:match}
For a trusted target context $x=(p,o,f,v)$---package, document origin, field class, and node token---$\Match(d,x)=1$ exactly when every mandatory component of $d$ accepts the corresponding component of $x$. An unavailable mandatory identity yields $0$, never a wildcard.
\end{definition}

\begin{definition}[Eligible handle]
\label{def:elig}
A record is \emph{eligible} at time $t$ for task $\tau'$ and target $x$ if its authenticator verifies, its state is live, $t_i\le t<t_e$, $u>0$, $\tau'=\tau$, $\Match(d,x)=1$, and the confirmation precondition $c$ is satisfied.
\end{definition}

\paragraph{Reference-monitor assumptions.} The guarantees below hold only under the following.

\begin{assumption}[Reference-monitor integrity]
\label{ass:tcb}
The adversary cannot modify the \SH{} service, store, resolver, redaction registry, trusted target-context extraction, injection path, or their control flow, and every resolve attempt passes through the resolver.
\end{assumption}

\begin{assumption}[Source and task authenticity]
\label{ass:source}
Only an approved source adapter can mint a protected-source handle, and the task token $\tau$ and provenance $\sigma$ cannot be forged or widened by the agent.
\end{assumption}

\begin{assumption}[Authenticator security]
\label{ass:crypto}
$F$ is existentially unforgeable under chosen-message attack, $K$ remains secret, identifiers are uniform over $\bits^\lambda$, and $\enc$ is canonical and injective over valid metadata.
\end{assumption}

\begin{assumption}[Authentic target identity]
\label{ass:identity}
For every accepted node token, the resolver-derived context $x=(p,o,f,v)$ equals the actual injection target at commit time; a required origin or field component is never treated as matched when it is unavailable.
\end{assumption}

\begin{assumption}[Modeled-channel redaction]
\label{ass:redaction}
Before $s$ becomes observable in a target field an obligation is active, and every screenshot, accessibility tree, event, action echo, and log included in the agent view applies the same value-independent mask until $s$ is absent.
\end{assumption}

\begin{assumption}[Value-oblivious public behavior]
\label{ass:oblivious}
For equal-length $s_0,s_1$ with equal declared leakage, minting, policy decisions, error messages, public timing buckets, and handle distributions are identically distributed.
\end{assumption}

\paragraph{Guarantees.} We now state the four properties that answer the attacker goals of Table~\ref{tab:goals}.

\begin{theorem}[Destination confinement]
\label{thm:confinement}
Under Assumptions~\ref{ass:tcb} and~\ref{ass:identity}, if a \SH{} value is committed to target $x$ under task $\tau'$, then $\Match(d,x)=1$ and $\tau'=\tau$ for the record bound to the handle.
\end{theorem}
\begin{proof}
A committed injection requires a prior successful reservation of the same record and target token. The resolver is the complete mediator (Assumption~\ref{ass:tcb}) and its transition guard is exactly eligibility (Definition~\ref{def:elig}), so reservation implies target and task acceptance. The injection capability is one-shot and binds the reservation to the revalidated node token, and Assumption~\ref{ass:identity} excludes target substitution before commit. Hence the committed target satisfies the predicate and task equality.
\end{proof}

\begin{corollary}[Redirection resistance]
\label{cor:redirection}
Under Theorem~\ref{thm:confinement}, if an attacker target $x_A$ is outside $d$, an injected instruction that requests \Enter{} at $x_A$ cannot cause the resolver to commit the secret there.
\end{corollary}
\begin{proof}
A commit at $x_A$ would imply $\Match(d,x_A)=1$, contradicting the premise.
\end{proof}

\begin{theorem}[Handle unforgeability]
\label{thm:unforgeability}
Let an adversary make at most $Q$ online acceptance attempts and observe handles of its choice through authorized mint events. Under Assumptions~\ref{ass:tcb} and~\ref{ass:crypto}, the probability that it causes acceptance of handle metadata not emitted by \Mint{} is at most $\Adv^{\mathsf{euf}}_{F,\mathcal{B}}(\lambda)+Q/2^\lambda$.
\end{theorem}
\begin{proof}
Construct a reduction $\mathcal{B}$ that simulates the store and answers authenticator queries through its MAC oracle. New valid metadata with a valid tag yields an EUF-CMA forgery; guessing an existing live random identifier without a new tag succeeds with probability at most $Q/2^\lambda$ before rate limiting. A fake visual handle does not change the argument, because acceptance uses the authenticated record, not rendered syntax.
\end{proof}

\begin{proposition}[Single-use linearizability]
\label{prop:singleuse}
Assume the live-to-reserved transition is an atomic compare-and-swap and no rollback occurs after possible delivery. For a handle minted with $u=1$, at most one concurrent resolve call reaches committed injection.
\end{proposition}
\begin{proof}
All contenders observe one state variable, so exactly one linearizes the transition from live to reserved; the others see reserved or consumed and reject. Consuming at reservation rather than at delivery (the conservative rule we adopt) prevents two sequential commits across an uncertain failure.
\end{proof}

\begin{proposition}[Expiry monotonicity]
\label{prop:expiry}
If expiry is evaluated against a monotonic clock and expired records never return to live, a handle rejected because $t\ge t_e$ cannot become eligible through a wall-clock adjustment.
\end{proposition}
\begin{proof}
Monotonic time does not decrease and no transition returns an expired record to live, so the strict conjunct $t<t_e$ of Definition~\ref{def:elig} remains false.
\end{proof}

\paragraph{Non-disclosure.} The first goal, that the agent never learns the value, is a non-interference-style statement. Let $\View_A(\pi)$ be the agent-visible projection of execution trace $\pi$: screenshots and accessibility trees, handle metadata, action results, timing rounded to a public granularity, and application messages, but not TCB memory, the private input buffer, or protected-confirmation pixels. Let $\ell$ be the declared leakage, containing the type $q$, a coarse destination label, age and remaining-lifetime buckets, and success/failure status, but never $s$, a deterministic hash of a low-entropy OTP, or value-dependent error codes.

\begin{definition}[Agent-view non-disclosure]
\label{def:nd}
For equal-length secrets $s_0,s_1$ with the same leakage $\ell$ and any probabilistic polynomial-time adversary $\mathcal{A}$ controlling the agent, the distinguishing advantage is
\begin{equation}
\begin{aligned}
\Adv^{\mathsf{nd}}_{\SH,\mathcal{A}}(\lambda)=\bigl|&\Pr[\mathcal{A}(\View_A(\Real(s_0)))=1]\\
&-\Pr[\mathcal{A}(\View_A(\Real(s_1)))=1]\bigr|.
\end{aligned}
\label{eq:nd}
\end{equation}
Non-disclosure holds when this advantage is negligible for equal $\ell$.
\end{definition}

\begin{theorem}[Agent-view non-disclosure]
\label{thm:nd}
Under Assumptions~\ref{ass:tcb}, \ref{ass:source}, \ref{ass:crypto}, \ref{ass:redaction}, and~\ref{ass:oblivious}, and before any explicitly declassified value-dependent outcome, $\Adv^{\mathsf{nd}}_{\SH,\mathcal{A}}(\lambda)$ in Eq.~\eqref{eq:nd} is negligible for equal-leakage secrets.
\end{theorem}
\begin{proof}
Consider a hybrid trace. Replacing the mint view for $s_0$ with that for $s_1$ changes the distribution negligibly: the visible identifier is uniform and independent of $s$, the type and public label are equal by the leakage premise, and the authenticator covers metadata that does not encode $s$ (Assumptions~\ref{ass:crypto},\ref{ass:oblivious}). Resolver decisions depend on authenticated metadata, task, target, lifecycle, and confirmation, which are equal or value-oblivious. At injection the raw value moves only through TCB memory and the target's private input path, and Assumption~\ref{ass:redaction} replaces every modeled post-entry occurrence with the same mask. Public errors and timing do not depend on the value (Assumption~\ref{ass:oblivious}), so the two agent-view distributions are indistinguishable until an allowed value-dependent outcome is released.
\end{proof}

\begin{remark}[Legitimate-but-unwanted actions]
Theorem~\ref{thm:confinement} does not cover a target inside $d$. An injection may still induce an unwanted action within the same authorized package or origin, such as logging into an attacker's account. Task semantics, transaction confirmation, and richer form binding are separate defenses; we quantify this residual risk rather than claim it eliminated.
\end{remark}
